\chapter{Геометрія JET: від плазми до котушок}\label{ch:geo6}

Попередні розділи йшли від $\psi$ до форми. Тут --- навпаки: від вимог до форми машини. Приклад --- проєкт JET 1975~р.~\cite{JET1975}, за яким проєкт має 3D-реконструкцію~\cite{ProjJET}. Повна «анатомія» з механікою, камерою, трансформатором і портами --- у методичці~\cite[розд.~5]{Manual}. Тут лише те, що стосується геометрії, і код, що її рахує.

\section{Аспектне відношення як компроміс}\label{sec:geo6-aspect}

Струм котушок TF, потрібний для заданих $\Ip$ і $q$, зростає як $(R_0/a)^2$~\cite[розд.~5, (5.1)]{Manual}, тому магніт дешевшає зі зменшенням $A$. Знизу $A$ обмежує центральний стовп: крізь отвір тора треба провести внутрішні ноги котушок TF і осердя трансформатора, яке наводить струм. Проєкт знайшов мінімум $A\approx2{,}2$ і обрав $R_0/a=2{,}37$~\cite[розд.~5]{Manual}. Табл.~\ref{tab:geo4-shape} показує той самий порядок у DIII-D ($A=2{,}74$). Ціна малого $A$ --- більше захоплених частинок: для зовнішньої поверхні JET $\sqrt{2\epsilon/(1+\epsilon)}=0{,}77$ при $\epsilon=a/R_0=0{,}42$~\cite[розд.~5]{Manual}.

\begin{figure}[!tb]\centering
\includegraphics[width=0.84\textwidth]{jet_section.jpg}
\caption{Вертикальний переріз JET 1975 з розмірами: від осі назовні --- центральна колона ярма, стек PF-1, внутрішня нога TF, камера, плазма ($\psiN$), зовнішня нога TF, PF-2--4, вертикальні плечі. Власна 3D-реконструкція~\cite{ProjJET} (рисунок з методички~\cite[розд.~5]{Manual}); побудовано за даними~\cite{JET1975}.}\label{fig:geo6-section}
\end{figure}

\section{Котушка TF: D чистого натягу}\label{sec:geo6-coil}

Котушка TF розтягується магнітним тиском назовні. Кругла котушка при цьому згинається, а згин руйнує ізоляцію між витками. D-подібна котушка сприймає навантаження \emph{лише натягом}, без згину. Умова цього --- радіус кривини контуру, пропорційний $R$~\cite[розд.~5, (5.3)]{Manual}:
\begin{equation}\label{eq:geo6-D}
\rho(R)=\kD R,\qquad \kD=\tfrac12\ln\frac{R_2}{R_1},
\end{equation}
де $R_1$ --- верх прямої внутрішньої ноги, $R_2$ --- зовнішній екватор. Для осьової лінії обмотки нашої моделі $R_1=\SI{1,304}{м}$, $R_2=\SI{4,790}{м}$~\cite[розд.~5]{Manual}, звідки $\kD=0{,}651$: радіус кривини $\SI{0,85}{м}$ біля внутрішньої ноги і $\SI{3,12}{м}$ на зовнішньому екваторі. Контур круто загинається всередині і полого --- зовні; разом із прямою ногою це й дає літеру D. Як уже сказано в п.~\ref{sec:geo4-miller}, це \emph{інша} крива, ніж D Міллера для плазми (див. також VR6: проєкт свідомо не наводить замкненої формули кривої Princeton D~\cite{RegViz}). У коді котушку не будують з~\eqref{eq:geo6-D}: контур отвору оцифровано з креслення проєкту (\texttt{machines/jet1975/profiles/tf\_coil\_inner.csv}) і зсунуто на половину товщини ноги.

\section{Гофрування тороїдального поля}\label{sec:geo6-ripple}

32 окремі котушки порушують осесиметрію: $|\Bvec|$ стає періодичним по $\tor$ з періодом $2\pi/32$. Глибину гофрування рахують так: $\dripple=(B_{\max}-B_{\min})/(B_{\max}+B_{\min})$, де $B_{\max}$ --- у площині котушки, $B_{\min}$ --- посередині між котушками. \emph{Пастка:} проєкт JET означує $\eripple=2\dripple$~\cite[розд.~5]{Manual}. Код рахує обидва варіанти за законом Біо--Савара від 32 ниток уздовж осьової лінії обмотки:

\code{viz/src/tokviz/machine/ripple.py}{122}{137}{(гофрування за Біо--Саваром)}

\noindent Рядки 131--133: $|\Bvec|$ у площині котушки ($\tor_0$) і посередині між котушками ($\tor_0+\pi/32$) у тих самих точках $(R,Z)$. Рядки 135--137: обидва означення. Відношення не залежить від струму, бо поле лінійне за ним.

\begin{projmodel}[гофрування JET]
Запуск \texttt{geom\_tour.py} (41~МА, одна нитка на котушку): на краю плазми $R=\SI{4,21}{м}$ $\dripple=1{,}83\,\%$, тобто $\eripple=3{,}67\,\%$ проти проєктних 3,6\,\%; на осі $R_0$ $\dripple=3\cdot10^{-7}$; $B(R_0)=\SI{2,770}{Тл}$, як у проєкті (2,77). Похибку дає переважно оцифрування контуру котушки: зсув осьової лінії на $\pm\SI{2}{см}$ змінює $\eripple(\SI{4,21}{м})$ на $+14/-12\,\%$~\cite[розд.~5]{Manual}.
\end{projmodel}

Чому гофрування так круто зростає назовні? Грубо (поза корпусом): поле дискретних ниток з кроком $2\pi R_2/N$ на радіусі $R_2$ спадає всередину як $(R/R_2)^N$. Для $N=32$ оцінка $(4{,}21/4{,}79)^{32}=1{,}6\,\%$ близька до 1,83\,\% за Біо--Саваром. Плазма JET займає отвір майже до зовнішньої ноги, і за це доводиться платити гофруванням.

\section{Рівновага в межі JET}\label{sec:geo6-solovev}

Щоб побачити магнітні поверхні всередині D-межі (п.~\ref{sec:geo4-jet}), потрібна рівновага. Найпростіша точна --- рівновага Соловйова: при сталих $\mu_0p'$ і $FF'$ рівняння Ґреда--Шафранова лінійне. Його розв'язок --- частинний поліном плюс сім однорідних розв'язків~\cite{Cerfon2010}. Сім коефіцієнтів задають так, щоб межа $\psi=0$ проходила через три точки D (зовнішній і внутрішній екватор, верх) з правильним нахилом і кривинами Міллера:

\code{viz/src/tokviz/machine/equilibrium_analytic.py}{283}{303}{(сім умов Cerfon--Freidberg)}

\noindent Рядки 283--285: $\epsilon=a/R_0$, $\arcsin\delta$ і частинний розв'язок $u_p$ ($A$ --- параметр поділу струму між тиском і $FF'$). Рядки 288--290: кривини межі Міллера на зовнішньому екваторі, внутрішньому екваторі й у верхній точці. Рядок 291: ці три точки в безрозмірних координатах $x=R/R_0$, $y=Z/R_0$. Рядки 296--300: сім умов --- $u=0$ у трьох точках, $\pa u/\pa x=0$ у верхній точці (там дотична горизонтальна) і три умови кривини. Рядок 303: лінійна система $7\times7$.

Лишаються два вільні параметри. Параметр $A$ підбирають бісекцією так, щоб $\betap=0{,}9$ (як у виносках проєкту), а масштаб $\psi_0$ --- так, щоб $\Ip=\SI{3,8}{МА}$:

\code{viz/src/tokviz/machine/equilibrium_analytic.py}{490}{517}{(збирання рівноваги JET)}

\noindent Рядки 501--506: $\kappa=b/a$, бісекція по $A$ для заданої $\betap$, розв'язок форми і масштаб $\psi_0$ під $\Ip$. Рядки 507--509: $\psi$ розраховують на сітці $193\times257$ із запасом за межею плазми і загортають у той самий клас \texttt{Equilibrium}, що й кадри EFIT. Тому весь код розд.~\ref{ch:geo2}--\ref{ch:geo5} (поверхні, $q$, $\thetastar$) працює для JET без змін.

\begin{table}[!tb]\centering\small
\caption{Рівновага Соловйова в межі JET 1975 (запуск \texttt{geom\_tour.py}; ті самі числа --- у \texttt{data/bake\_jet1975/physics.json}~\cite{ProjJET})}\label{tab:geo6-eq}
\begin{tabular}{@{}ll@{\hspace{2.5em}}ll@{}}\toprule
$A$ (поділ струму) & 0,0948 & $\Ip$ & 3,80~МА \\
$\psi_0$ & 15,64~Вб/рад & $\betap$ / $\betat$ & 0,900 / 2,32\,\% \\
нев'язка межі & 0,08~мм & $\li(1)$ / $\li(3)$ & 0,498 / 0,425 \\
вісь $\Rax$ & 3,183~м & зсув Шафранова $\Rax-R_0$ & 0,223~м \\
$q_0$ (з гесіана) & 2,80 & $\qnf$ / $q(0{,}99)$ & 5,51 / 5,75 \\
об'єм (квадратура) & 145,46~м$^3$ & об'єм (Папп, табл.~\ref{tab:geo4-shape}) & 145,46~м$^3$ \\\bottomrule
\end{tabular}
\end{table}

\begin{figure}[!tb]\centering
\includegraphics[width=0.86\textwidth]{geo_jet_solovev.jpg}
\caption{Рівновага Соловйова в межі JET 1975. Ліворуч: поверхні $\psiN=0{,}1\ldots0{,}9$, межа Міллера ($\kappa=1{,}68$, $\delta=0{,}348$) і магнітна вісь, зсунута на \SI{0,22}{м} назовні від $R_0$. Праворуч: $q(\psiN)$ з контурного інтеграла~\eqref{eq:geo3-q} і $q_0$ з гесіана~\eqref{eq:geo5-q0} (точка). Власний розрахунок, код: \texttt{geometry/code/geom\_tour.py}.}\label{fig:geo6-solovev}
\end{figure}

Об'єм, обчислений квадратурою по сітці, збігається з об'ємом за Паппом (п.~\ref{sec:geo4-params}) з відносною різницею $10^{-5}$ --- незалежна перевірка обох кодів. Контурний $q$ і $q$ з трасування силових ліній збігаються до 0,015\,\% на $\psiN=0{,}1\ldots0{,}95$ (\texttt{physics.json}, ключ \texttt{q\_traced\_vs\_contour}). Помітний і зсув Шафранова: вісь на \SI{0,22}{м} зовні від $R_0$ (рис.~\ref{fig:geo6-solovev}), а поверхні згущуються до зовнішнього боку.

\emph{Чесне обмеження.} Струм Соловйова пласкої форми: $\Jtor=(C_1R+C_2/R)/\mu_0$ від $\psi$ не залежить. Тому $q_0=2{,}80$ і $q(a)/q_0\approx2{,}1$, тоді як проєкт розраховував на пікований струм з $q_0\approx1$. Поверхонь $q=1$, $3/2$, $2$ у цій моделі немає. Для топології (вкладені поверхні, замикання раціональних ліній $q=3,\,7/2,\,4,\,5$ з точністю $<\SI{0,5}{мм}$) модель придатна, для МГД-стійкості --- ні~\cite[розд.~5]{Manual}.

\begin{practicum}[JET]
(1)~Побудуйте рівновагу з \texttt{method="lsq"} замість \texttt{"cf7"}: як зміняться нев'язка межі і $q_0$? (2)~Змініть $\delta$ на 0 і на 0,5 (\texttt{build\_equilibrium(delta=...)}): як залежать $\qnf$ і зсув Шафранова від трикутності? (3)~Порахуйте $\dripple$ на $R=4{,}0\ldots4{,}4$~м з кроком 5~см і порівняйте з оцінкою $(R/R_2)^{32}$. Де оцінка найгірша і чому?
\end{practicum}
